\section{Introduction}

Le fonctionnement d'un véhicule moderne repose sur une électronique complexe qui coordonne de
multiples calculateurs embarqués, ou \emph{Electronic Control Units} (\ecus). Ces calculateurs
forment ce que l'on nomme des \emph{systèmes réactifs}~: ils interagissent en permanence avec le
monde physique et doivent réagir instantanément et correctement à des événements externes, comme une
pression sur la pédale de frein ou un signal radar. Lorsque la défaillance d'un tel système peut
avoir des conséquences graves, il est qualifié de \emph{critique}. Cette thèse CIFRE, réalisée au
sein d'\ampere, la division logicielle du Groupe Renault, aborde la conception des logiciels pour
ces systèmes réactifs critiques embarqués.

\subsection{Transition d'Architecture Logicielle}

L'architecture électronique traditionnelle des véhicules repose sur un réseau distribué d'\ecus
spécialisés. Chaque fonctionnalité, ou presque, possède son propre calculateur~: l'un pour la
gestion moteur, l'autre pour les airbags. Cependant, l'intégration de fonctionnalités avancées,
telles que l'assistance à la conduite (ADAS) ou l'intelligence artificielle, rend ce modèle
décentralisé inadapté, difficile à faire évoluer et à maintenir.

Face à cette complexité croissante, l'industrie automobile opère une transition majeure vers une
architecture centralisée~: le \emph{Software-Defined Vehicle} (\sdv). Ce paradigme remplace de
nombreux \ecus par un nombre restreint de calculateurs haute performance. Cette centralisation
simplifie l'architecture matérielle, réduit le câblage et facilite les mises à jour logicielles à
distance. Cette évolution, bien que nécessaire, introduit deux défis majeurs pour l'ingénierie
logicielle.

\subsection{Les Défis de l'Architecture \sdv}

\paragraph{Défi n°1~: Changement du mode de communication}
Le premier défi concerne l'inefficacité du mode de communication hérité des architectures
distribuées. Dans ce modèle, les composants communiquent de manière \emph{périodique}, en diffusant
leur état à intervalle régulier, même en l'absence de nouvelle information. Si cette approche est
adaptée aux bus traditionnels de communication, elle se révèle problématique pour les architectures
\sdv qui exploitent des réseaux à haut débit comme Ethernet.
%
Le bus CAN, historiquement utilisé, est conçu pour être efficace avec de petits volumes de données~:
son protocole est léger et sa surcharge minimale. À l'inverse, Ethernet repose sur des protocoles
réseaux standards (TCP/IP) qui encapsulent les données à transmettre dans plusieurs couches. Chaque
couche ajoute ses propres en-têtes, créant une surcharge protocolaire significative~: transmettre
quelques octets utiles peut nécessiter l'envoi d'un paquet de plusieurs dizaines d'octets.

Une expérience menée chez \ampere a mis en lumière cette problématique. La configuration du système,
basée sur une transmission périodique de tous les signaux à 10~ms, a provoqué une saturation
complète du bus de communication. L'analyse a révélé que pour certains calculateurs, 97\% des
messages transmis n'étaient pas utilisés (soit superflus, soit perdus). Ce trafic inutile a non
seulement gaspillé la bande passante et les ressources de calcul, mais a également engendré des
pertes de données.
%
Cet exemple démontre qu'un changement d'approche est indispensable. Un modèle de communication
\emph{événementiel}, où les données ne sont transmises que lorsqu'une information pertinente est
produite, serait plus approprié. Cette approche exige cependant une refonte des modèles de
programmation traditionnels.

\paragraph{Défi n°2~: Fossé d'abstraction des outils de programmation}
Le second défi réside dans l'adoption du langage de programmation \Rust. Bien que plébiscité par
l'industrie pour ses garanties de sûreté mémoire et sa performance, \Rust reste un langage de
\emph{bas niveau} qui impose aux développeurs une gestion manuelle de nombreux détails
d'implémentation.

Or, l'histoire de l'ingénierie a montré que la programmation bas niveau des systèmes critiques est
une source de risques notoire. Le cas tristement célèbre de l'appareil de radiothérapie Therac-25
en est une illustration dramatique~: des erreurs logicielles subtiles (conditions de concurrence,
dépassements mémoire) dans un système programmé en assembleur ont conduit à des surdoses mortelles
de radiations pour plusieurs patients~\cite{journals/computer/LevesonT93}.

À cette complexité s'ajoute le choix d'\ampere d'utiliser la programmation \emph{asynchrone} en
\Rust pour répondre au défi des communications événementielles. Ce modèle, bien qu'efficace, est
réputé pour sa difficulté et sa propension à introduire des erreurs d'un nouveau genre, comme les
interblocages. Cette double complexité---bas niveau et asynchrone---crée un \emph{fossé
      d'abstraction} entre les outils de modélisation de haut niveau (tels que
\Simulink~\cite{simulink}), familiers aux ingénieurs système, et les exigences de l'implémentation.
Demander aux experts du contrôle-commande de maîtriser ces détails est une source d'erreurs et un
frein à la productivité.

\subsection{Contribution}

Pour répondre à ces défis, cette thèse introduit \GRust, un langage de programmation dédié
(\textit{Domain Specific Language}, ou \dsl) qui vise à combler ce fossé d'abstraction. \GRust
permet aux ingénieurs de spécifier le comportement des systèmes réactifs à un haut niveau
d'abstraction, et de générer le code \Rust asynchrone, sûr et performant, adapté aux architectures
\sdv.

Pour ce faire, \GRust s'appuie sur le principe \emph{Globalement Asynchrone, Localement Synchrone}
(GALS)~\cite{DBLP:phd/us/Chapiro85} et combine deux paradigmes~:
\begin{itemize}
      \item Un modèle \emph{synchrone}, inspiré des langages comme \Lustre~\cite{%
                  DBLP:conf/ifip/Kahn74,%
                  DBLP:conf/popl/CaspiPHP87,%
                  DBLP:conf/tase/ColacoPP17%
            }. Ce paradigme repose sur l'hypothèse où les calculs sont considérés comme instantanés.
            Le système progresse par une succession d'instants logiques discrets, garantissant un
            comportement déterministe. Cette logique de calcul est décrite dans \GRust \via le
            sous-langage \GRSync, offrant un cadre familier, conceptuellement similaire aux blocs de
            \Simulink.

      \item Un modèle \emph{asynchrone} et événementiel, inspiré de la Programmation Fonctionnelle
            Réactive (\frp)~\cite{%
                  DBLP:conf/icfp/ElliottH97,%
                  DBLP:conf/haskell/Elliott09,%
                  DBLP:conf/padl/WanTH02%
            }. Ce modèle, implémenté dans le sous-langage \GRFRP, orchestre les flux de données
            entre les blocs de calcul, assurant que les interactions ne se produisent que lorsque
            cela est nécessaire.
\end{itemize}

De plus, \GRust permet de spécifier des contraintes temporelles pour garantir la sûreté. Par
exemple, une fréquence minimale d'émission assure la \emph{vivacité} des composants (prouvant qu'ils
sont toujours actifs), tandis qu'une fréquence maximale prévient la surcharge du système.

Au-delà de la conception du langage, notre contribution inclut un \emph{compilateur} qui traduit les
spécifications \GRust en code \Rust asynchrone, masquant la complexité inhérente à ce paradigme.
Elle inclut également une formalisation sémantique du langage. Cette formalisation, à travers deux
sémantiques---opérationnelle et dénotationnelle---dont l'équivalence est prouvée, garantit des
propriétés essentielles comme la productivité (un programme ne se bloque jamais) et le déterminisme.
Nous introduisons aussi un théorème de sur-échantillonnage, qui valide la correction de
l'interaction entre les mondes synchrone et asynchrone.
%
Enfin, pour les systèmes les plus critiques, la vérification formelle est impérative. À cette fin,
nous avons développé \GReusot, un outil qui fait le lien avec le vérificateur de programmes \Rust
\Creusot~\cite{DBLP:conf/icfem/DenisJM22}. \GReusot permet d'annoter les modèles \GRust avec des
propriétés, rendant la vérification formelle accessible au niveau du modèle, sans nécessiter une
expertise du code généré.

\medskip

En conclusion, cette thèse propose une solution qui permet à l'industrie automobile de tirer parti
des architectures \sdv et du langage \Rust. \GRust comble le fossé d'abstraction en offrant un
modèle de conception de haut niveau, tout en générant un code sûr, efficace et formellement
vérifiable, sans compromettre les méthodologies de développement établies.

\section{Description Technique de \GRust}

\GRust est un langage conçu pour combler le fossé entre la modélisation de haut niveau, familière
aux ingénieurs système, et la programmation bas niveau en \Rust asynchrone, requise par les
architectures \sdv.

\subsection{Description Syntaxique}

Cette section décrit les éléments syntaxiques et les concepts fondamentaux du langage \GRust. Nous y
présentons, à l'aide de courts exemples, les quatre composantes qui structurent le langage~: le
noyau asynchrone \GRFRP pour la définition de \emph{service}, la bibliothèque d'opérateurs \GReact,
le noyau synchrone \GRSync pour la logique de calcul, et l'extension \GReusot pour la vérification
formelle.

\subsubsection{\GRFRP~: Le Noyau Asynchrone pour les Services}

\GRFRP est le langage de haut niveau qui permet de définir l'architecture globale du système sous
la forme de \emph{services}. Un service orchestre les flux de données asynchrones entre les
interfaces du système (capteurs, actuateurs) et la logique de calcul interne.

Les flux de données sont de deux types. Les \emph{signaux} (\fkf{signal}) représentent des
valeurs continues dans le temps, comme la vitesse du véhicule. Les \emph{événements}
(\fkf{event}) représentent des occurrences discrètes et ponctuelles, comme un clic de bouton.

L'interface d'un \kft{service} est déclarée \via les mots-clés \kft{import} et \kft{export}. La
logique interne est ensuite définie de manière déclarative dans le corps du \kft{service}.
Les contraintes temporelles sont spécifiées par l'annotation \kft{@}\interval{10}{3000}, indiquant
que le service \texttt{adaptive\_cruise\_control} ne calcule pas plus souvent que toutes les 10~ms
et au moins toutes les 3~s.

\begin{lstlisting}[language=GRust]
// Interface du service
import signal speed_km_h: float;
export signal brakes_m_s: float;

// Définition du service avec contraintes temporelles (min 10ms, max 3s)
service adaptive_cruise_control @[10, 3000] {
    // Déclaration d'un flux local
    let signal speed_m_s: float = convert(speed_km_h);
    // Définition d'un flux exporté
    brakes_m_s = acc(speed_m_s, ...);
}
\end{lstlisting}

\subsubsection{\GReact~: La Bibliothèque Réactive Intégrée}

Pour manipuler ces flux asynchrones, \GRFRP est accompagné de \GReact, une `boîte à outils'
d'opérateurs intégrée, inspirée de \textsc{ReactiveX}~\cite{reactivex}. Elle fournit des fonctions
de flux prêtes à l'emploi pour les scénarios réactifs courants. En voici quelques éléments:

\begin{description}
      \item[Changement de type] L'opérateur \texttt{on\_change} permet de créer un évènement à
            partir d'un signal, ne contenant que ses changement de valeur.
            \begin{lstlisting}[language=GRust, numbers=none]
let event radar_event: float = on_change(radar_signal);
\end{lstlisting}
      \item[Gérer le temps] L'opérateur \texttt{timeout} déclenche un événement si un évènement
            reste silencieux trop longtemps, utile pour la surveillance.
            \begin{lstlisting}[language=GRust, numbers=none]
let event watchdog_event: unit = timeout(radar_event, 1000); // 1s
\end{lstlisting}
      \item[Fusionner des évènements] L'opérateur \texttt{merge} combine deux événements en un seul.
            \begin{lstlisting}[language=GRust, numbers=none]
let event any_user_action: unit = merge(button_press, screen_touch);
\end{lstlisting}
\end{description}

\subsubsection{\GRSync~: Le Noyau Synchrone pour les Composants}

\GRSync est le langage dédié à la création d'opérateurs de flux. Il est synchrone et déclaratif,
inspiré de \Lustre~\cite{DBLP:conf/popl/CaspiPHP87}. On peut y définir~:
\begin{itemize}
      \item Des \emph{fonctions}, qui sont des calculs purs, sans état.
      \item Des \emph{composants}, qui peuvent maintenir un état interne entre les instants de
            calcul.
\end{itemize}

L'état est géré \via l'opérateur \kft{last}, qui accède à la valeur précédente d'un signal. Pour la
logique événementielle, \GRSync utilise deux constructions principales~:

\begin{itemize}
      \item \kft{when}~: Permet de mettre a jour des signaux et d'envouer des évènements en réaction
            à des occurrences discrètes. Un bloc \kft{init} définit la valeur initiale des signaux,
            puis d'autres branches sont activées par la présence d'un événement ou d'un front
            montant (un signal booléen qui passe de \kft{false} à \kft{true}). Les signaux non mis à
            jour conservent leur valeur, créant une mémoire implicite, et les évènements sont
            absents.
            \begin{lstlisting}[language=GRust]
component counter(max: int, count: unit?) -> (c: int, reset: unit?) {
      when {
            init    => { c = 0; }
            count?  => { c = 1 + last c; }
            c > max => { c = 0; reset = emit (); }
      }
}
\end{lstlisting}
      \item \kft{match}~: Structure de contrôle classique pour sélectionner un mode de calcul en
            fonction de la valeur d'un signal, à la manière d'un \kft{switch} en \Lustre.
            \begin{lstlisting}[language=GRust]
match mode {
      Mode::Eco   => { consumption = compute_eco(...); }
      Mode::Sport => { consumption = compute_sport(...); }
}
\end{lstlisting}
\end{itemize}

\subsubsection{\GReusot~: Vérification Formelle Intégrée}

Pour garantir la sûreté exigée par la norme \iso, \GRust met à disposition l'extension \GReusot qui
permet d'intégrer des \emph{contrats formels} directement dans le code \GRSync.
Il est possible de spécifier des \emph{préconditions} (\kft{requires}) et des \emph{postconditions}
(\kft{ensures}) pour les fonctions et composants, leur donnant une spécification de preuve formelle,
en plus de la spécification calculatoire.

\begin{lstlisting}[language=GRust]
function safe_divide(a: int, b: int) -> int
    // Précondition : le diviseur ne doit pas être nul
    requires { b != 0 }
{
    return a / b;
}
\end{lstlisting}
Le compilateur traduit ces annotations dans le code \Rust pour le vérificateur
\Creusot~\cite{DBLP:conf/icfem/DenisJM22}, qui peut alors prouver mathématiquement que le programme
respecte ses contrats (par exemple, qu'une division par zéro n'arrivera jamais). Cette approche rend
la vérification formelle plus accessible en la maintenant au niveau du modèle.

\subsection{Vérifications Statiques à la Compilation}

Pour garantir la robustesse et la prévisibilité des programmes, le compilateur \GRust réalise une
série de vérifications statiques avant de générer le code \Rust. Ces analyses permettent de détecter
de nombreuses erreurs au plus tôt, directement au niveau du modèle.

\paragraph{Gestion des identifiants et de leur portée} Le compilateur vérifie que chaque identifiant
(variable, composant, \etc) est déclaré avant d'être utilisé et qu'il n'y a pas de redéfinition
ambiguë. Il applique pour cela deux modèles de portée distincts~:
\begin{description}
      \item[Portée lexicale] Dans un service et une fonction, les identifiants sont définis de
            manière séquentielle. Un identifiant local n'est visible que pour les instructions qui
            le suivent.
      \item[Portée récursive] Dans un composant, toutes les équations sont mutuellement récursives.
            Cela permet de décrire des systèmes de dataflow où les dépendances
            sont cycliques (par exemple, \texttt{a = 1 + last b; b = a;}), ce qui est courant en
            contrôle-commande.
\end{description}
De plus, des règles strictes s'appliquent aux structures de contrôle. Par exemple, toutes les
branches d'un \kft{match} doivent définir exactement le même ensemble d'identifiants, assurant
qu'une sortie est toujours calculée. À l'inverse, un \kft{when} permet des mises à jour
partielles, ce qui correspond à sa sémantique événementielle.

\paragraph{Système de typage} \GRust implémente son propre vérificateur de types pour fournir des
messages d'erreur clairs et spécifiques au domaine des flots de données. Ce système est fondamental
pour faire le pont entre les deux mondes du langage~:
\begin{itemize}
      \item Dans \GRFRP, on manipule des \emph{types de flux}~: \signal{\texttt{T}} et
            \event{\texttt{T}}, qui décrivent le comportement temporel d'une donnée ayant des
            valeurs de type \texttt{T}.
      \item Dans \GRSync, on manipule des \emph{types échantillonnés}~: un \signal{\texttt{T}} est
            échantillonné en une valeur de type \texttt{T} (car toujours présente), tandis qu'un
            \event{\texttt{T}} est échantillonné en une valeur de type \opt{\texttt{T}}
            (car possiblement absente).
\end{itemize}
Cette distinction permet au compilateur de rejeter des constructions sémantiquement incorrectes,
comme l'application d'un opérateur temporel tel que \kft{timeout} à un signal, ou l'utilisation
d'un événement dans une opération arithmétique sans avoir préalablement résolu sa présence.
L'obligation de noter explicitement tous les types renforce la lisibilité et la prévisibilité du
code.

\paragraph{Analyse de causalité} Une troisième vérification statique de \GRust, hérité des langages
synchrones, est la \emph{causalité}. Un programme est dit causal si chacune de ses sorties peut être
calculée en un nombre fini d'étapes à partir des entrées. Cette propriété garantit qu'un programme
est productif et ne contient pas de dépendances cycliques instantanées (par exemple,
\texttt{x = x + 1}), qui créeraient un paradoxe logique insoluble. Le compilateur \GRust intègre une
passe d'\emph{analyse de causalité} pour vérifier statiquement cette propriété sur les composants.

\subparagraph{Dépendances et Ordre de Calcul} L'analyse de causalité repose sur la construction d'un
\emph{graphe de dépendances} pour chaque composant. Les nœuds du graphe représentent les
identifiants (entrées, sorties, variables locales), et les arêtes représentent les dépendances de
données entre eux. Par exemple, l'équation \texttt{x = y + z} crée des arêtes de \texttt{y} et
\texttt{z} vers \texttt{x}. Une dépendance \via un opérateur \kft{last} est marquée comme
`retardée', car elle dépend d'une valeur calculée à l'instant précédent.
%
Un composant est causal si et seulement si son graphe de dépendances ne contient aucun cycle de
dépendances instantanées. Un tel cycle signifierait qu'une valeur dépend d'elle-même au même instant
de calcul, ce qui est impossible à résoudre.

\subparagraph{Gestion des Dépendances Conditionnelles} La complexité de l'analyse augmente avec les
constructions de contrôle comme \kft{match}. Une branche peut introduire une dépendance qui
n'existe pas dans une autre. Par exemple~:
\begin{lstlisting}[language=GRust, numbers=none]
match c { true  => { x = y; } false => { y = x; } }
\end{lstlisting}
Une analyse naïve conclurait à une dépendance cyclique entre \texttt{x} et \texttt{y}. Pour éviter
ce problème, le compilateur utilise une technique de \emph{renommage}~\cite{bourke:hal-03936656}~:
les variables définies dans chaque branche sont traitées comme des entités distinctes (par exemple,
\texttt{x\_true}, \texttt{y\_false}). Le graphe de dépendances est construit sur ces variables
renommées, ce qui permet de capturer précisément les flux de données conditionnels sans créer de
faux cycles.

\subparagraph{Ordonnancement des Équations} Si un composant est causal, cela garantit qu'il existe
au moins un ordre d'évaluation valide (un ordonnancement) pour ses équations. Le compilateur utilise
le graphe de dépendances pour effectuer un \emph{tri topologique} des opérations, produisant ainsi
la séquence d'instructions impératives nécessaire pour traduire un composant en une machine à états.

\subsection{Implémentation du Compilateur}

Le prototype de compilateur \GRust est implémenté en \emph{macro procédurale} \Rust, appelée
\texttt{grust!}. Plutôt qu'un programme externe, le compilateur s'intègre directement dans la chaîne
de compilation \Rust standard (\texttt{cargo build}). Il transforme une stream de token
correspondant au code \GRust en stream de token correspondant au code \Rust asynchrone. Cette
approche délègue la responsabilité de la génération du code machine final au compilateur \Rust, et
permet une interopérabilité avec le code \Rust existant.

\paragraph{Compilation de \GRSync} Un composant \GRSync, défini par des équations déclaratives, est
compilé en une machine à états impérative. La cible est une structure \Rust qui implémente un trait
\texttt{Component}, exposant deux méthodes~: \texttt{init} pour créer l'état initial, et
\texttt{step} pour exécuter un pas de calcul. Pour cela, le compilateur effectue plusieurs passes~:
\begin{description}
      \item[Analyse grammaticale] Un \emph{Arbre de Syntaxe Abstrait} est créé à partir du code
            source. Il represente le programme de manière structurée, en accord avec la grammaire de
            \GRust.
      \item[Analyses statiques] S'en suit les analyses présentées précédemment.
      \item[Normalisation] Les constructions de haut niveau comme \kft{when} sont traduites en
            expressions \kft{match} complexes, unifiant la logique de contrôle.
      \item[Ordonnancement] En utilisant l'ordre topologique issu de l'analyse de causalité, les
            équations sont ordonnancées en une séquence d'instructions linéaires.
      \item[Génération de code] La séquence ordonnancée devient le corps de la méthode
            \texttt{step}, et les variables d'état (issues des opérateurs \kft{last}) deviennent
            les champs de la structure de la machine à états.
\end{description}

\paragraph{Compilation de \GRFRP} Un service \GRFRP est compilé en une \emph{machine de propagation
      asynchrone}. La cible est une fonction \texttt{run} qui lance une tâche asynchrone gérant le
cycle de vie du service. Le principe de fonctionnement est le suivant~:
\begin{enumerate}
      \item Toutes les sources d'entrée (flux importés, timers internes) sont fusionnées en un
            unique flux, qui est ensuite trié par ordre de priorité. Pour l'instant l'ordre est
            chronologique, mais il est possible de définir un ordre plus complexe prenant en compte
            la criticité des différentes entrées.
      \item Une boucle principale asynchrone attend le prochain événement de ce flux ordonné.
      \item Lorsqu'un événement arrive, il est distribué à un gestionnaire (\emph{handler}) dédié,
            généré spécifiquement pour cette entrée.
      \item Ce gestionnaire exécute la chaîne de calculs impactée par cette entrée, en appelant les
            méthodes \texttt{step} des composants concernés. Si une sortie est modifiée, la nouvelle
            valeur est envoyée de manière asynchrone vers le bus de communication.
\end{enumerate}
Ce modèle garantit une exécution efficace (les calculs ne sont effectués que si nécessaire) et
déterministe (les événements sont traités dans un ordre déterminé statiquement).

\paragraph{Compilation de \GReusot} Le compilateur traduit les contrats \GReusot en annotations
\Creusot sur le code \Rust généré. Pour une fonction, il génère également une version
\emph{logique} parallèle, opérant sur des types mathématiques (\eg entiers naturels), et ajoute une
preuve que le code exécutable se comporte conformément à ce modèle logique. Cela permet une
vérification compositionnelle où la preuve d'un composant peut s'appuyer sur les spécifications
des fonctions qu'il appelle.

\subsection{Évaluations et Benchmarks}

L'évaluation de \GRust, menée en partie lors de deux stages (l'un académique, l'un industriel),
s'articule en trois volets pour valider ses contributions principales~: la performance de son noyau
synchrone, son expressivité pour la conception de haut niveau, et son applicabilité dans un
contexte industriel \sdv.

\subsubsection{Évaluation des performances de \GRSync}

Le premier volet évalue les performances de \GRSync, en le comparant aux compilateurs académiques
\Lustre de référence (\Heptagon~\cite{heptagon}, \LustreC~\cite{lustrec} et \Velus~\cite{velus}).
Pour ce faire, un frontend, \GLustre, a été développé pour transpiler du \Lustre en \GRSync.
Plusieurs systèmes réalistes, comme le filtre de Kalman du drone \crazyflie~\cite{crazyflie-fm} et
un modèle d'avion de ligne (TCM)~\cite{tcmOriginal,tcmCocospec}, ont été utilisés pour les tests.
%
Les résultats montrent que la chaîne d'outils de \GRSync est très compétitive par rapport aux outils
de l'état de l'art académique.

\subsubsection{Évaluation de la conception de haut niveau}

Le deuxième volet est une étude de cas qualitative, centrée sur le même filtre de Kalman du drone
\crazyflie, pour comparer l'expressivité de \GRust à une implémentation bas niveau. L'implémentation
originelle en \Ci repose sur des mécanismes de synchronisation manuels et complexes (sémaphores)
pour gérer la concurrence entre la tâche d'estimation et la boucle de contrôle principale. Le code
résultant est difficile à analyser et source potentielle d'erreurs.

À l'inverse, le modèle \GRust remplace cette complexité par un modèle déclaratif et événementiel où
les dépendances de données sont explicites. La communication entre composants se fait \via des flux
typés, éliminant le besoin de synchronisation manuelle. Le modèle \GRust est ainsi significativement
plus concis, lisible et, par construction, plus sûr.

\subsubsection{Comparaison avec du code \Rust industriel}

Le troisième volet, mené lors d'un stage industriel chez \ampere, compare le code généré par \GRust
à des implémentations manuelles en \Rust asynchrone pour trois services automobiles réels. Ces
services ont été réimplémentés en \GRust et déployés sur le calculateur cible du véhicule pour une
mesure en conditions réelles. Deux scénarios d'entrées ont été utilisés~: un flux \emph{périodique}
(données envoyées à intervalles fixes) et un flux \emph{événementiel} (données envoyées uniquement
en cas de changement).
%
Les résultats sont les suivants~:
\begin{itemize}
      \item \textbf{Concision~:} Les modèles \GRust sont jusqu'à 71\% plus concis que le code
            manuel, démontrant un gain de productivité et de maintenabilité significatif.

      \item \textbf{Validation fonctionnelle~:} La comparaison des sorties a montré une équivalence
            fonctionnelle entre le code généré et le code manuel.

      \item \textbf{Performance~:} Le passage d'une communication périodique à événementielle divise
            la charge CPU par deux, validant l'approche promue pour les architectures \sdv. Le code
            généré par \GRust affiche des performances (CPU, mémoire) quasi identiques à celles du
            code manuel optimisé, avec un surcoût négligeable.

      \item \textbf{Analyse de traces~:} L'analyse des traces d'exécution (\via
            \perfetto~\cite{perfetto}) confirme que le gain en performance provient de l'élimination
            des traitements redondants sur les entrées inchangées.

      \item \textbf{Validation du théorème de sur-échantillonnage~:} En comparant les sorties de
            \GRust entre les scénarios périodique et événementiel, nous avons validé
            expérimentalement le théorème de sur-échantillonnage.
\end{itemize}

En conclusion, ces évaluations démontrent que \GRust atteint ses objectifs. Il offre un niveau
d'abstraction élevé qui améliore la productivité et la sûreté, tout en générant un code dont les
performances sont comparables à celles d'un développement manuel bas niveau, ce qui en fait une
solution viable et attractive pour l'industrie automobile.

\subsection{Synthèse}

Cette section a détaillé la conception et l'implémentation de \GRust, de sa syntaxe à son évaluation
en passant par les mécanismes de son compilateur. Les benchmarks ont montré que le code généré est
non seulement fonctionnellement correct, mais aussi très performant, rivalisant les compilateurs
académiques de référence et égalant les performances de code \Rust écrit à la main par des experts.

\section{Sémantique Formelle}

Cette section établit les fondations théoriques de \GRust en définissant une sémantique formelle
pour ses deux noyaux, \GRSync et \GRFRP. Pour chaque sous-langage, deux modèles sémantiques
complémentaires sont développés~: une \emph{sémantique opérationnelle} qui décrit l'exécution pas à
pas, et une \emph{sémantique dénotationnelle} offre une vision sur une exécution complète.
L'équivalence de ces deux modèles est prouvée, tout comme les propriétés de productivité et de
déterminisme. Une contribution centrale est la résolution formelle du problème de
\emph{sur-échantillonnage} \via un système de types réactifs, garantissant la correction sémantique
de l'interaction entre le monde asynchrone des services et le monde synchrone des composants.

\subsection{Sémantique Opérationnelle~: Modéliser l'Exécution Pas à Pas}

La sémantique opérationnelle décrit comment un programme \GRust évolue en réponse à des entrées.
Elle modélise le comportement du système tel qu'il serait exécuté.

\subsubsection{Composant \GRSync~: Machine à États}

Le comportement d'un composant \GRSync modélise une \emph{machine à états finis}. La sémantique
définit une fonction de transition qui, pour un état interne donné et un échantillon des entrées à
un instant logique, produit un nouvel état et les sorties correspondantes.

\paragraph{État et Transition d'un Composant} L'état d'un composant encapsule sa mémoire~: la
valeur précédente de chaque signal accédé \via \kft{last}, ainsi que l'état interne de tous les
sous-composants qu'il instancie. La sémantique d'un composant est donc une fonction de transition
qui prend l'état courant et les entrées actuelles pour calculer le nouvel état et les sorties.

\paragraph{Évaluation des Équations par Point Fixe} Le corps d'un composant est un ensemble
d'équations simultanées avec des dépendances mutuelles. Pour modéliser leur évaluation en un seul
instant logique, la sémantique utilise une \emph{itération par point fixe}.
\begin{enumerate}
      \item Au début du calcul, tous les identifiants définis localement sont inconnus et marqués
            comme indéfinis ($\bot$).
      \item Le système évalue itérativement l'ensemble des équations. À chaque itération, les
            équations dont les dépendances sont maintenant connues sont calculées.
      \item Si une valeur n'est pas encore calculée, toute opération qui en dépend produit $\bot$.
      \item Le processus continue jusqu'à ce qu'un point fixe soit atteint, c'est-à-dire que toutes
            les valeurs sont calculées.
\end{enumerate}
L'analyse de causalité (décrite à la section précédente) garantit que ce point fixe est toujours
atteint en un nombre fini d'itérations.

\subsubsection{Service \GRFRP~: Machine de Propagation}

Un service \GRFRP modélise une \emph{machine de propagation asynchrone}. Son rôle est de réagir à
des messages externes et de propager les changements à travers les dépendances des flux.

\paragraph{État et Messages} L'état d'un service mémorise la dernière valeur connue des signaux et
des évènements dont ses opérations dépendent. Le service réagit à des messages entrants, qui peuvent
être~: un changement de valeur sur un flux importé, muni d'un horodatage (\textit{timestamp}) ; ou
l'expiration d'un minuteur interne, également horodaté.

\paragraph{Propagation des Changements} Lorsqu'un message est reçu, la sémantique décrit comment
cette information déclenche une cascade de calculs~:
\begin{enumerate}
      \item Le message initial met à jour le contexte de propagation pour le flux correspondant.
      \item Le système parcourt ensuite les dépendances du graphe de flux. Pour chaque opérateur, il
            évalue si ses entrées dans le contexte ont changé.
      \item Si une entrée a changé, l'opérateur est ré-évalué, ce qui peut produire de nouvelles
            valeurs. Ces nouvelles valeurs mettent à leur tour à jour le contexte, continuant la
            propagation.
      \item Si un flux exporté est mis à jour, un message de sortie est généré.
\end{enumerate}
Ce modèle décrit une exécution où un calcul n'est effectué qu'en réponse à un changement.

\paragraph{Opérateurs \GReact} La sémantique formalise le comportement de chaque opérateur. Par
exemple, \timeoutExpr{\texttt{e}}{\texttt{T}} est modélisé par deux règles, une qui réinitialise
un minuteur interne lorsque l'événement \texttt{e} arrive, et une autre qui produit un événement de
sortie lorsque le minuteur correspondant expire (le minuteur est réinitialisé également).

\paragraph{Interaction avec \GRSync} Le lien entre les deux mondes est formalisé au niveau de
l'appel de composants dans un service. Lorsqu'un service \GRFRP propage un changement vers un
composant \GRSync, il déclenche la fonction de transition \texttt{step} du composant. Cette
transition n'a lieu que si au moins une des entrées du composant a effectivement changé. Les
résultats de ce calcul synchrone sont ensuite propagés à leur tour dans le service.

\subsection{Sémantique Dénotationnelle~: Une Vision Abstraite des Flux}

La sémantique dénotationnelle s'intéresse à la signification d'un programme sur une exécution
infinie. Chaque flux est interprété comme une suite de valeurs infinie, appelée \emph{stream}. Ce
formalisme repose sur la \emph{coinduction}~\cite{DBLP:journals/mscs/KozenS17}, un principe de
preuve adapté aux structures de données infinies.

\subsubsection{Composant \GRSync~: Fonction de Streams Synchrones}

Dans le monde synchrone de \GRSync, le temps progresse par \emph{instants logiques} discrets. Un
flux est donc une séquence infinie de valeurs où le $n^\text{ième}$ élément correspond à la valeur
du flux au $n^\text{ième}$ instant. Pour gérer les calculs qui ne sont pas actifs à chaque instant
(par exemple, dans une branche de \kft{match}), on utilise une valeur spéciale $\bot$ (bottom) pour
marquer l'absence de valeur, garantissant ainsi que tous les flux restent temporellement alignés.

Un composant est interprété comme une fonction de stream, qui transforme des streams d'entrée en
streams de sortie. Les équations du composant sont vues comme des contraintes d'égalité sur ces
streams.
\begin{itemize}
      \item L'opérateur \lastExpr{\texttt{x}} est modélisé par une fonction
            \buffer{\texttt{c}}{\texttt{s}}, qui décale la stream \texttt{s} correspondant à
            \texttt{x} d'un instant et le préfixe avec une valeur initiale.

      \item Les constructions \kft{match} et \kft{when} sont interprétées \via un \emph{calcul
                  d'horloges} implicite. Une horloge est unz stream de booléens qui dicte quand un
            calcul est actif. Chaque branche d'un \kft{match} ou d'un \kft{when} définit une
            \emph{sous-horloge} qui n'est active que lorsque la condition de la branche est vraie.
            Les calculs de la branche sont alors effectués sur cette sous-horloge. Pour un
            \kft{match}, l'exhaustivité des patterns garantit que la fusion des sous-horloges de
            chaque branche recompose l'horloge d'origine. Pour un \kft{when}, la sémantique est
            différente~: lorsqu'aucune branche n'est active, les signaux conservent leur valeur
            précédente (\via un opérateur \kft{last} implicite) et les événements sont absents.
\end{itemize}
Ce modèle, classique dans la théorie des langages synchrones, permet de raisonner sur le
comportement temporel global du système de manière purement fonctionnelle.

\subsubsection{Service \GRFRP~: Fonction de Stream Horodatées}

Dans le monde asynchrone de \GRFRP, un flux est interprété comme une stream de valeurs horodatées
(\emph{timestamped stream}), c'est-à-dire une séquence de paires \texttt{(t, v)} où les horodatages
sont strictement croissants. Ce modèle, inspiré du Tagged-Signal
Model~\cite{DBLP:journals/tcad/LeeS98}, capture la nature non-synchronisée des communications.

Un service est interprété comme une fonction qui transforme un contexte de stream d'entrée
horodatées en un contexte de stream de sortie horodatées. Les opérateurs de la bibliothèque \GReact
sont définis comme des fonctions coinductives sur ces streams~:
\begin{itemize}
      \item \timeoutExpr{\texttt{s}}{\texttt{T}} est une fonction qui compare les horodatages de la
            stream d'entrée \texttt{s} à une deadline interne pour générer des événements de sortie.
      \item \mergeExpr{\texttt{s1}}{\texttt{s2}} est une fonction qui fusionne deux streams en une
            seule de telle sorte que l'horodatage reste croissant, en priorisant la stream de
            gauche \texttt{s1}.
\end{itemize}
La validité de ces fonctions coinductives est prouvée, garantissant la correction de la sémantique.

\paragraph{Le Pont entre \GRFRP et \GRSync} La sémantique dénotationnelle formalise également
l'interaction entre les deux noyaux. Lorsqu'un service \GRFRP appelle un composant \GRSync, un pont
sémantique est établi~:
\begin{enumerate}
      \item Les streams d'entrée horodatées du service sont projetées sur une horloge synchrone~:
            pour un signal, sa dernière valeur est répétée à chaque instant logique ; pour un
            événement, sa valeur est présente uniquement à l'instant de son occurrence et absente
            ($\bot$) sinon.
      \item La fonction de streams synchrone du composant \GRSync est ensuite appliquée.
      \item Les streams synchrones de sortie sont re-basculées dans le monde asynchrone pour devenir
            des streams horodatées, ne conservant que les changements ou occurrences.
\end{enumerate}
Cette double projection assure que la sémantique synchrone et déterministe des composants est
correctement intégrée dans le cadre global asynchrone des services.

\subsection{Théorèmes Fondamentaux de Correction}

Sur la base des sémantiques opérationnelle et dénotationnelle, nous prouvons des théorèmes qui
garantissent la prévisibilité des programmes \GRust. Ces preuves sont menées séparément pour les
deux noyaux, \GRSync et \GRFRP, en respectant leurs sémantiques distinctes.

\subsubsection{Productivité}

Le théorème de productivité garantit qu'un programme \GRust passant les vérifications (typage,
identifiants, causalité) ne se bloque jamais et produira toujours une sortie pour toute entrée
valide (en l'absence d'erreurs d'exécution). Pour un système réactif qui doit fonctionner en
continu, c'est une propriété fondamentale qui assure sa vivacité.

Le théorème énonce que la fonction de transition de la sémantique opérationnelle est totale. Pour
tout état valide et toute entrée valide, la fonction produit toujours une sortie et un nouvel état
valides. Pour le prouver, on définit une fonction d'observation qui exécute le programme sur une
séquence d'entrées et on démontre que cette fonction est toujours définie.
%
La preuve est menée par induction structurelle sur les constructions du langage.
\begin{itemize}
      \item Pour les constructions simples (constantes, opérations), la preuve est directe.
      \item Pour les constructions complexes, la preuve s'appuie sur l'hypothèse d'induction
            appliquée à leurs sous-expressions.
      \item La principale difficulté, les équations mutuellement récursives de \GRSync, est résolue
            par une \emph{induction sur le rang de causalité} des identifiants. Le rang d'un
            identifiant correspond à la longueur de sa plus longue chaîne de dépendances. En
            prouvant la propriété pour les identifiants de rang 0 (sans dépendances), puis en
            montrant qu'elle se propage aux rangs supérieurs, on garantit que chaque valeur peut
            être calculée en un nombre fini d'itérations.
\end{itemize}

\subsubsection{Déterminisme}

Le théorème de déterminisme assure que le comportement d'un programme est prévisible et
reproductible. Il est indispensable pour le test, le débogage et la certification, car il garantit
que les sorties dépendent uniquement des entrées, et non de facteurs externes non spécifiés comme
l'ordonnancement des tâches.

Le théorème énonce que pour un historique d'entrées donné, l'historique des sorties est unique.
Formellement, il est prouvé au niveau de la sémantique dénotationnelle~: si deux exécutions partent
de streams d'entrée équivalentes, alors leurs streams de sortie sont également équivalentes.
%
La preuve procède également par induction structurelle sur le langage, en montrant que chaque
opérateur est une fonction mathématique.
\begin{itemize}
      \item Pour les opérateurs simples, la preuve est directe.
      \item Pour les opérateurs de \GReact, on montre que les fonctions coinductives qui les
            définissent sont déterministes.
      \item Pour les équations récursives de \GRSync, on utilise à nouveau une induction sur le rang
            de causalité pour montrer que chaque stream est uniquement déterminée par les streams
            dont elle dépend.
\end{itemize}

\subsubsection{Équivalence}

Le théorème d'équivalence prouve que la sémantique opérationnelle (le modèle d'exécution concret)
est une implémentation fidèle de la sémantique dénotationnelle (la spécification mathématique
abstraite). Cette cohérence garantit que les propriétés prouvées sur le modèle abstrait s'appliquent
bien au comportement réel du programme.

Le théorème établit qu'une simulation infinie de la sémantique opérationnelle sur une séquence
d'entrées produit une trace de valeurs qui est une dénotation valide, et inversement.
%
La preuve est menée par coinduction, la technique appropriée pour raisonner sur les stream infinis.
\begin{itemize}
      \item La direction \emph{opérationnelle $\Rightarrow$ dénotationnelle} est la plus complexe.
            Elle se fait par induction structurelle sur le programme, en montrant que chaque pas de
            la sémantique opérationnelle respecte les contraintes d'égalité de la sémantique
            dénotationnelle.

      \item La direction \emph{dénotationnelle $\Rightarrow$ opérationnelle} est plus simple. Elle
            s'appuie sur les théorèmes de productivité et de déterminisme. On sait qu'il existe
            \emph{au moins une} exécution opérationnelle (grâce à la productivité). On montre que
            cette exécution satisfait la sémantique dénotationnelle (grâce à la première direction
            de la preuve). Enfin, comme la sémantique dénotationnelle est déterministe, on conclut
            que \emph{toutes} les exécutions possibles sont équivalentes, y compris celle que l'on
            cherchait à construire.
\end{itemize}

\subsection{Le Défi du Sur-échantillonnage et sa Solution}

Une contribution théorique majeure de la thèse est la résolution du conflit sémantique entre
l'exécution événementielle des services \GRFRP et la logique synchrone des composants \GRSync. Cette
section détaille le problème, la solution apportée par le système de types réactifs, et la
validation formelle \via le théorème de sur-échantillonnage.

\subsubsection{Le Problème du Sur-échantillonnage (\textit{Oversampling})}

L'architecture GALS de \GRust implique qu'un composant synchrone \GRSync est exécuté par un service
asynchrone \GRFRP. Idéalement, pour être efficace, le service ne devrait déclencher un pas de calcul
du composant que lorsque ses entrées changent. Cependant, un service peut être activé par une entrée
qui n'est pas une dépendance du composant, et il est possible que le composant soit tout de même
ré-évalué~: il est \emph{sur-échantillonné}.

Ce sur-échantillonnage devient problématique si la logique du composant dépend de flux qui ne sont
pas stables entre deux de ses propres activations. Il existe trois sources principales d'instabilité~:
\begin{description}
      \item[Le temps physique] L'opérateur \timeExpr de \GRFRP fournit l'horodatage de l'événement
            déclencheur. Si un composant dépend de cette valeur, chaque sur-échantillonnage (qui a
            un nouvel horodatage) modifiera son état, même si ses entrées fonctionnelles n'ont pas
            changé.
      \item[Les mémoires non gardées] Un compteur comme \texttt{o = last o + 1} a une boucle de
            rétroaction interne non conditionnée. Chaque exécution, y compris un
            sur-échantillonnage, incrémentera son état.
      \item[Les émissions d'événements non gardées] Un appel non gardé de \emitExpr{\texttt{x}}
            produit un évènement qui est sémantiquement toujours présent. Une exécution
            événementielle ne le déclenchera cependant que lors d'un changement de \texttt{x},
            créant une incohérence.
\end{description}
Dans ces cas, le comportement du composant dépendrait de la fréquence d'activation du service, et
non plus seulement de ses propres entrées, violant ainsi le principe de compositionnalité.

\subsubsection{La Solution~: un Système de Types Réactifs}

Pour résoudre ce problème, un \emph{système de types réactifs} est introduit. Son rôle est de
garantir statiquement qu'un sur-échantillonnage n'aura pas d'effets de bord non désirés. Pour ce
faire, il classifie tous les flux en deux catégories~:
\begin{description}
      \item [Discrets] Ce sont les flux dont la valeur ne change qu'en réponse à une modification
            explicite d'une de leurs dépendances directes. Ils sont stables lors d'un
            sur-échantillonnage.
      \item [Continus] Ce sont les flux dont la valeur peut changer lors d'un sur-échantillonnage,
            dont les trois sources d'instabilité~: \kft{time}, les mémoires \kft{last} et les
            émissions \kft{emit} non gardées.
\end{description}

Le système de types impose que les mémoires \kft{last} et la logique événementielle (les patterns
d'un \kft{when}) ne peuvent dépendre que de flux discrets. Cette règle, appliquée à la fois dans
\GRSync et \GRFRP, empêche statiquement les dépendances temporelles implicites qui causent les
problèmes de sur-échantillonnage. La construction \kft{when} devient le mécanisme central de
\emph{discrétisation}, en s'assurant que les calculs ne sont activés que par des événements discrets.

\subsubsection{Le Théorème de Sur-échantillonnage}

Ce typage est validé par le \emph{Théorème de Sur-échantillonnage}. Pour tout service ou composant
bien typé comme étant réactif, une évaluation due au sur-échantillonnage (c'est-à-dire une exécution
alors que ses entrées discrètes n'ont pas changé) ne modifie ni son état interne, ni la valeur de
ses sorties discrètes. Ce théorème prouve que le modèle d'exécution événementiel de \GRust est
sémantiquement correct. Il justifie que l'on peut opter pour une exécution évènementielle sans
sacrifier la prévisibilité et le déterminisme du comportement synchrone.

La preuve est menée par induction structurelle sur les constructions du langage, pour \GRSync et
\GRFRP.
\begin{itemize}
      \item Pour \GRSync, on formalise un pas de sur-échantillonnage en fournissant au composant des
            entrées qui sont des \emph{non-changements} (la même valeur pour un signal, ou l'absence
            pour un événement). On prouve alors que l'état interne reste identique et que les
            sorties discrètes sont également des \emph{non-changements}. La principale difficulté
            est de montrer que les constructions de contrôle (\kft{match}, \kft{when}) se comportent
            de manière stable, ce qui est garanti par le système de types réactifs qui assure que
            leurs conditions de garde sont discrètes.

      \item Pour \GRFRP, le sur-échantillonnage est modélisé par l'arrivée d'un message sur un flux
            que le service n'importe pas. La preuve montre que cette arrivée ne déclenche aucune
            chaîne de propagation pour les flux discrets. Pour raisonner sur la stabilité des
            valeurs à travers plusieurs pas de temps, une fonction auxiliaire
            \texttt{current\_value} est introduite pour calculer la valeur courante d'un flux à
            partir de son historique complet de changement. On démontre alors que la valeur courante
            d'une sortie discrète après une séquence d'exécutions est la même, que cette séquence
            contienne ou non des pas de sur-échantillonnage.
\end{itemize}
Ensemble, ces preuves établissent une fondation solide pour l'exécution événementielle de \GRust.

\subsection{Synthèse}

Cette section a établi les fondations théoriques de \GRust en développant un cadre formel pour ses
deux noyaux, \GRSync et \GRFRP. Pour chacun, nous avons défini deux sémantiques~:
\begin{itemize}
      \item La \emph{sémantique opérationnelle} modélise l'exécution pas à pas du langage. Elle
            décrit les composants \GRSync comme des machines à états finis et les services \GRFRP
            comme des machines de propagation asynchrone.

      \item La \emph{sémantique dénotationnelle} fournit une interprétation mathématique abstraite,
            associant à chaque flux une \textit{stream}. Elle modélise les composants comme des
            fonctions sur les stream et les équations comme des contraintes, permettant de raisonner
            sur le comportement global du programme.
\end{itemize}
La \emph{productivité} garantit qu'un programme ne se bloque jamais. Le \emph{déterminisme} assure
que son comportement est prévisible et reproductible. Enfin, le théorème d'\emph{équivalence} prouve
que les deux sémantiques coïncident, assurant que le modèle d'exécution concret est une
implémentation fidèle de la spécification mathématique.

Une contribution théorique majeure de cette section a été la résolution formelle du problème de
\emph{sur-échantillonnage}. Nous avons montré comment une exécution événementielle pouvait entrer
en conflit avec la sémantique synchrone, et nous avons résolu ce conflit en introduisant un
\emph{système de types réactifs}. Ce système distingue les flux \emph{continus} (sensibles au
temps d'exécution) des flux \emph{discrets} et impose des règles strictes pour leur interaction.
La validité de cette approche a été confirmée par le \emph{théorème de sur-échantillonnage}, qui
prouve qu'un programme bien typé est sémantiquement insensible aux pas de calcul superflus.

\section{Conclusion et Perspectives}

Cette thèse a abordé les défis logiciels critiques qui émergent de la transition de l'industrie
automobile vers \emph{Software-Defined Vehicle} (\sdv). Pour répondre à ces enjeux, nous avons
introduit \GRust, un langage de programmation dédié qui réconcilie les pratiques de modélisation
établies avec les exigences des architectures logicielles modernes. \GRust unifie la rigueur
déterministe du paradigme \emph{synchrone}, familier aux ingénieurs, avec la flexibilité de la
\emph{Programmation Fonctionnelle Réactive} (\frp) pour modéliser les comportements événementiels.
Cette synthèse permet de concevoir des systèmes complexes à un haut niveau d'abstraction, tout en
générant automatiquement un code \Rust asynchrone, sûr et performant.

Les contributions de cette thèse sont les suivantes~:
\begin{enumerate}
      \item la conception du langage \GRust et de son compilateur prototype ;

      \item l'établissement des fondations théoriques du langage à travers des sémantiques
            opérationnelle et dénotationnelle ;

      \item la productivité, le déterminisme et l'équivalence des deux sémantiques ;

      \item la résolution du conflit sémantique entre l'exécution événementielle des services et la
            logique synchrone des composants, que nous nommons sur-échantillonnage, résolu par un
            système de types réactifs ; et

      \item la validation pratique de l'approche par des benchmarks.
\end{enumerate}

\subsection{Limites des Travaux Actuels}

Bien que cette thèse établisse la viabilité de l'approche \GRust, certaines limites définissent le
périmètre des contributions actuelles. Le compilateur prototype implémente le système de types
réactifs, mais aucune aide compréhensible n'est pour l'instant fournie pour corriger un programme ne
respectant pas ces règles. Les preuves sémantiques ont été réalisées manuellement et n'ont pas
été vérifiées par un assistant de preuve formel comme \Rocq~\cite{rocq} ou
Isabelle~\cite{DBLP:books/sp/NipkowPW02}, ce qui représenterait un niveau de confiance supérieur.
Enfin, l'outil de vérification \GReusot est limité dans sa capacité à raisonner de manière
compositionnelle sur des composants à état, car il ne permet pas encore de spécifier ou d'utiliser
des invariants d'état pour les sous-composants.

\subsection{Perspectives et Travaux Futurs}

Les résultats prometteurs de \GRust ouvrent de nombreuses perspectives de recherche et de
développement pour en faire un outil industriel robuste.

\paragraph{Évolution du Langage et du Compilateur} La priorité est d'ajouter des aides à la
correction d'un programme ne respectant pas le système de types réactifs, sur la base d'heuristiques
sur les cas d'erreurs fréquents. À plus long terme, le compilateur pourrait être rendu agnostique au
\textit{runtime} asynchrone (actuellement dépendant de \texttt{tokio}~\cite{tokio}) pour cibler des
environnements embarqués plus contraints.

\paragraph{Vérification et Validation} Pour renforcer les garanties de sûreté, l'outil \GReusot doit
être étendu pour supporter la vérification compositionnelle de composants à état, notamment par
l'ajout d'invariants. De plus, la création d'un framework de test intégré à \GRust permettrait aux
ingénieurs de valider leurs modèles par rapport à des spécifications fonctionnelles. Enfin, la
vérification formelle du \textit{runtime} asynchrone lui-même et l'intégration d'outils d'analyse de
temps d'exécution au pire cas (WCET) compléteraient la chaîne de preuves, de la conception à
l'exécution.

\paragraph{Industrialisation et Écosystème} Pour une adoption industrielle, la chaîne d'outils
\GRust doit être enrichie pour faciliter la qualification selon la norme \iso. Cela inclut la
génération automatique d'artefacts de certification, comme des matrices de traçabilité et des
documents de conception. Des benchmarks sur des cas d'usage plus complexes, comme des systèmes ADAS
complets, seront également nécessaires pour valider la scalabilité de l'approche.

\subsection{Remarques Finales}

En conclusion, cette thèse a démontré qu'une approche par modélisation de haut niveau, incarnée par
\GRust, est une stratégie viable et efficace pour maîtriser la complexité croissante des logiciels
automobiles. En unifiant la rigueur du paradigme synchrone avec l'expressivité de la programmation
réactive, \GRust comble le fossé entre les outils de conception traditionnels et les exigences des
architectures logicielles modernes. En générant un code \Rust sûr et performant, il offre une voie
pragmatique vers le développement de la prochaine génération de systèmes critiques, en garantissant
à la fois la sûreté, la maintenabilité et la performance.
